Appending CST coefficients as unknowns to a coupled viscous/inviscid Newton
solver, with geometric design constraints expressed as the linear rows CST's
own linearity makes available, converts constrained inverse airfoil design
from a search over an objective function into a single square, determined
root-find. The falsifiable self-consistency test of
Section~\ref{sec:falsifiable} passes with a wide margin on every accuracy and
iteration-count criterion, including a release-and-verify check the original
protocol had not required, though its convergence rate is fast superlinear
rather than quadratic (Section~\ref{sec:h3}).
The ablation matrix of Section~\ref{sec:results} confirms the architecture's
viability across three NACA sections and Bernstein orders up to a conditioning
cliff at $n=12$, and the two pre-registered panels extend the recovery claim
to 18 of 18 eligible NACA sections and 83 of 83 converged UIUC sections,
subject to the selection effect Section~\ref{sec:h1-uiuc-panel} quantifies.

Two of the study's results correct the hypothesis stated in
Section~\ref{sec:headline}. The cost
advantage over a competently-tuned nested baseline is $3.1$ to $8.1\x$ in
counted flow solves and $3.4$ to $3.5\x$ in wall-clock time, not the two to three
orders of magnitude assumed a priori. And the convergence behaviour is fast
superlinear rather than guaranteed quadratic, limited by the
finite-difference flow-block Jacobian and, under interpolated station
addressing, by a neglected term in the target rows. Against those corrections
stands a methodological contribution the original hypothesis did not
anticipate: QR-pivoted station selection as the identifiability guard for
square inverse systems, without which the same solver converges cleanly to a
wrong root while reporting success. The paper's central results are therefore
a falsified numeric hypothesis and a strengthened, unanticipated qualitative
one.

\paragraph{Future work.} The single most valuable next experiment is the one
Section~\ref{sec:validity} names: recovering a geometry from a target this
method did not generate, whether from a finer discretisation, a perturbed
distribution, or a published pressure distribution, together with a sweep of
recovery error against target noise. Until that is run, the precision reported
here is a statement about identifiability rather than about accuracy.

Beyond it, two extensions are pre-planned and neither is yet built. Cascade
periodicity substitutes a periodic kernel for the isolated-airfoil
panel-influence kernel, replacing $\ln(z-z_0)$ with $\ln[\sin(\pi(z-z_0)/s)]$
for pitch $s$, which reduces to the isolated-airfoil kernel as $s\to\infty$
and therefore comes with a free regression test; it would be validated against
exact cascade solutions before any turbine-cascade inverse-design claim is
made. Carrying the architecture into MISES itself uses that solver's modal
inverse mode. Because CST's perturbation modes are identical to its own basis
functions, a direct consequence of the linearity of
Section~\ref{sec:cst-linearity}, this requires writing $C(\psi)S_i(\psi)$ into
the solver's mode-file format with no modification to its source, and would
demonstrate the hypothesis in a transonic, viscous, real-cascade solver at
near-zero implementation cost if that route succeeds.

The gap between what this paper can currently claim and what was originally
pre-registered is threefold. The single-run ablation evidence has not been
extended to panel level. Three UIUC sections still fail the host solver's
wake-direction assertion for a cause unrelated to sharp trailing edges. And
the constraint rows, which are half of the formulation's novelty claim, remain
implemented and tested but unexercised by any measured result.
